\section{AHSME 1990}
        \begin{enumerate}
        	\item (\textit{AHSME 1990}) If $\dfrac{\frac{x}{4}}{2}=\dfrac{4}{\frac{x}{2}}$, then $x=$
$\textbf{(A)}\ \pm\frac{1}{2}\qquad\textbf{(B)}\ \pm 1\qquad\textbf{(C)}\ \pm 2\qquad\textbf{(D)}\ \pm 4\qquad\textbf{(E)}\ \pm 8$

 \textbf{Answer: }E



\textbf{Solution}

Cross-multiplying leaves 

 
 \begin{align*}\dfrac{x^2}{8} &= 8\\ x^2 &= 64\\ \sqrt{x^2} &= \sqrt{64}\\ x &= \pm 8\end{align*}


 So the answer is $\boxed{\text{(E)} \, \pm 8}$.

 


	\item (\textit{AHSME 1990}) $\left(\frac{1}{4}\right)^{-\frac{1}{4}}=$
$\textbf{(A) } -16\qquad \textbf{(B) } -\sqrt{2}\qquad \textbf{(C) } -\frac{1}{16}\qquad \textbf{(D) } \frac{1}{256}\qquad \textbf{(E) } \sqrt{2}$

 \textbf{Answer: }E



\textbf{Solution}

$\left(\frac14\right)^{-\frac14}=4^{\frac14}=\sqrt[4]4=\sqrt{\sqrt4}=\sqrt2$ which is $\fbox{E}$



	\item (\textit{AHSME 1990}) The consecutive angles of a trapezoid form an arithmetic sequence. If the smallest angle is $75^\circ$, then the largest angle is


 $\textbf{(A) } 95^\circ\qquad \textbf{(B) } 100^\circ\qquad \textbf{(C) } 105^\circ\qquad \textbf{(D) } 110^\circ\qquad \textbf{(E) } 115^\circ$

 \textbf{Answer: }C



\textbf{Solution 1}

Since a trapezoid is a quadrilateral, its interior angles sum to $360^\circ$.

 Let $d$ be the common difference of the arithmetic sequence.  Then the trapezoid's angle measures, in degrees, are $75, 75+d, 75+2d, 75+3d$, so $300+6d=360$ and $d=10$. Thus the largest angle is $75 + 3 \cdot 10 = \boxed{\textbf{(C)}\ 105^\circ}$.

 



\textbf{Solution 2}

Like in all quadrilaterals, the interior angles of a trapezoid sum to $360^\circ$.

 Recall that the sum of an arithmetic sequence with $n$ terms, first term $a_1$, and last term $a_n$ is $\frac{n}{2}(a_1 + a_n)$.  Substituting $n = 4$ and $a_1 = 75^\circ$ yields $\frac{4}{2}(75^\circ + a_4) = 360^\circ$, and solving this equation yields $a_4 = \boxed{\textbf{(C)}\ 105^\circ}$.

 


	\item (\textit{AHSME 1990}) 
                    \begin{figure}[H]
                        \centering
                        \includegraphics[width=0.5\textwidth]{Problem/AHSME_I_1990_4_Problem_Figure_1.png}
                    \end{figure}
                    


 Let $ABCD$ be a parallelogram with $\angle{ABC}=120^\circ, AB=16$ and $BC=10.$ Extend $\overline{CD}$ through $D$ to $E$ so that $DE=4.$ If $\overline{BE}$ intersects $\overline{AD}$ at $F$, then $FD$ is closest to


 $\textbf{(A) } 1\qquad \textbf{(B) } 2\qquad \textbf{(C) } 3\qquad \textbf{(D) } 4\qquad \textbf{(E) } 5$

 \textbf{Answer: }B



\textbf{Solution}

$DFE$ and $AFB$ are similar triangles, so $FD$ is one quarter the length of the corresponding side $AF$.Thus it is one fifth of the length of $AD$, which means $\fbox{B}$



	\item (\textit{AHSME 1990}) Which of these numbers is largest?


 $\textbf{(A) } \sqrt{\sqrt[3]{5\cdot 6}}\qquad \textbf{(B) } \sqrt{6\sqrt[3]{5}}\qquad \textbf{(C) } \sqrt{5\sqrt[3]{6}}\qquad \textbf{(D) } \sqrt[3]{5\sqrt{6}}\qquad \textbf{(E) } \sqrt[3]{6\sqrt{5}}$

 \textbf{Answer: }B



\textbf{Solution}

Putting them all in the form $\sqrt{\sqrt[3]N}$ (the order of the radicals doesn't matter), the resulting interior numbers are 
 \[30, 1080, 750, 150, 180\] so the answer is $\fbox{B}$



	\item (\textit{AHSME 1990}) Points $A$ and $B$ are $5$ units apart. How many lines in a given plane containing $A$ and $B$ are $2$ units from $A$ and $3$ units from $B$?


 $\textbf{(A) } 0\qquad \textbf{(B) } 1\qquad \textbf{(C) } 2\qquad \textbf{(D) } 3\qquad \textbf{(E) }\text{more than }3$

 \textbf{Answer: }D



\textbf{Solution}

The lines have to be tangent to both of these circles.
                    \begin{figure}[H]
                        \centering
                        \includegraphics[width=0.5\textwidth]{Solution/AHSME_I_1990_6_Solution_Figure_1.png}
                    \end{figure}
                    By inspection, $\fbox{D}$



	\item (\textit{AHSME 1990}) A triangle with integral sides has perimeter $8$. The area of the triangle is


 $\textbf{(A) } 2\sqrt{2}\qquad \textbf{(B) } \frac{16}{9}\sqrt{3}\qquad \textbf{(C) } 2\sqrt{3} \qquad \textbf{(D) } 4\qquad \textbf{(E) } 4\sqrt{2}$

 \textbf{Answer: }A



\textbf{Solution}

The shortest side must be $1$ or $2$. However none of $(1,1,6),(1,2,5),(1,3,4)$ form triangles, so the shortest side must be $2$. Then $(2,2,4)$ is degenerate, so the sides must be $(2,3,3)$.

 This can be cut in half and reassembled into a rectangle with one side $1$ and diagonal $3$. By Pythagoras its area is then $2\sqrt2$ which means $\fbox{A}$



	\item (\textit{AHSME 1990}) The number of real solutions of the equation

 \[|x-2|+|x-3|=1\]
is


 $\textbf{(A) } 0\qquad \textbf{(B) } 1\qquad \textbf{(C) } 2\qquad \textbf{(D) } 3\qquad \textbf{(E) } \text{more than } 3$

 \textbf{Answer: }E



\textbf{Solution}

For $2\le x\le 3$, the left-hand side is $(x-2)-(x-3)$ which is $1$ for all $x$ in the interval. $\fbox{E}$



	\item (\textit{AHSME 1990}) Each edge of a cube is colored either red or black. Every face of the cube has at least one black edge. The smallest number possible of black edges is


 $\textbf{(A) } 2\qquad \textbf{(B) } 3\qquad \textbf{(C) } 4\qquad \textbf{(D) } 5\qquad \textbf{(E) } 6$

 \textbf{Answer: }B



\textbf{Solution}

Each black edge can only take care of two adjoining faces, so we know at least three will be needed. Once the first black edge is placed, it is easy to see that three will be sufficient, if they are separated and go in different directions:
                    \begin{figure}[H]
                        \centering
                        \includegraphics[width=0.5\textwidth]{Solution/AHSME_I_1990_9_Solution_Figure_1.png}
                    \end{figure}
                    This gives the answer $\fbox{B}$



	\item (\textit{AHSME 1990}) An $11\times 11\times 11$ wooden cube is formed by gluing together $11^3$ unit cubes. What is the greatest number of unit cubes that can be seen from a single point?


 $\textbf{(A) }  328\qquad \textbf{(B) }  329\qquad \textbf{(C) }  330\qquad \textbf{(D) }  331\qquad \textbf{(E) }  332$

 \textbf{Answer: }D



\textbf{Solution 1}

It is impossible to see two opposite sides of the cube from a single point.  Therefore, at most three faces of the cube may be seen from a single point, and these three faces must share a vertex.  There does exist a point where three such faces are visible.  Let $P$ be such a point.

 The large cube is made of $11^3 = 1331$ unit cubes.  The unseen cubes form a $10\times 10 \times 10$ cube, so $10^3=1000$ cubes are not visible from $P$.  Therefore, $1331-1000=\boxed{\textbf{(D)}\ 331}$ cubes are visible from $P$.

 


 



\textbf{Solution 2}


                    \begin{figure}[H]
                        \centering
                        \includegraphics[width=0.5\textwidth]{Solution/AHSME_I_1990_10_Solution_Figure_1.png}
                    \end{figure}
                    The best angle for cube viewing is centered on the corner, meaning three of the six faces are visible. So, therefore, the answer is just counting the number of cubes on the three faces, which is $3 \times 10^2$ for the inside parts of the faces, plus $3 \times 10$ for the edges, plus $1$ for the single shared cube in the corner, giving a total of $331$ or $\fbox{D}.$



	\item (\textit{AHSME 1990}) How many positive integers less than $50$ have an odd number of positive integer divisors?


 $\textbf{(A) } 3\qquad \textbf{(B) } 5\qquad \textbf{(C) } 7\qquad \textbf{(D) } 9\qquad \textbf{(E) } 11$

 \textbf{Answer: }C



\textbf{Solution}

Divisors come in pairs, unless there is an integer square root, so we just need the perfect squares below $50$. There are $7$, so $\fbox{C}$



	\item (\textit{AHSME 1990}) Let $f$ be the function defined by $f(x)=ax^2-\sqrt{2}$ for some positive $a$. If $f(f(\sqrt{2}))=-\sqrt{2}$ then $a=$
$\textbf{(A) } \frac{2-\sqrt{2}}{2}\qquad \textbf{(B) } \frac{1}{2}\qquad \textbf{(C) } 2-\sqrt{2}\qquad \textbf{(D) } \frac{\sqrt{2}}{2}\qquad \textbf{(E) } \frac{2+\sqrt{2}}{2}$

 \textbf{Answer: }D



\textbf{Solution}

If $f(x)=-\sqrt2$, then $ax^2=0\implies x=0$. Therefore $f(\sqrt2)=0\implies 2a=\sqrt2$, so $\fbox{D}$



	\item (\textit{AHSME 1990}) If the following instructions are carried out by a computer, which value of $X$ will be printed because of instruction $5$?


 <pre>1. START $X$ AT $3$ AND $S$ AT $0$.  
2. INCREASE THE VALUE OF $X$ BY $2$.   
3. INCREASE THE VALUE OF $S$ BY THE VALUE OF $X$. 
4. IF $S$ IS AT LEAST $10000$,   
       THEN GO TO INSTRUCTION $5$;  
       OTHERWISE, GO TO INSTRUCTION $2$.  
       AND PROCEED FROM THERE.  
5. PRINT THE VALUE OF $X$.  
6.STOP. 
</pre>$\textbf{(A) } 19\qquad \textbf{(B) } 21\qquad \textbf{(C) } 23\qquad \textbf{(D) } 199\qquad \textbf{(E) } 201$

 \textbf{Answer: }E



\textbf{Solution}

Looking at the first few values, it becomes clear that the program stops when 
 \[5+7+9+11+13+\ldots+(x-2)+x\ge 10000\] which is to say 
 \[1+3+5+7+9+11+13+\ldots+(x-2)+x\ge 10004\]However, the left hand side is now simply the square $\frac{(x+1)^2}4$. Multiplying out, we get 
 \[x+1\ge \sqrt{40016}\approx 200.039996...\]

 So the correct answer is $201$, which is $\fbox{E}$



	\item (\textit{AHSME 1990}) 
                    \begin{figure}[H]
                        \centering
                        \includegraphics[width=0.5\textwidth]{Problem/AHSME_I_1990_14_Problem_Figure_1.png}
                    \end{figure}
                    


 An acute isosceles triangle, $ABC$, is inscribed in a circle. Through $B$ and $C$, tangents to the circle are drawn, meeting at point $D$. If $\angle{ABC}=\angle{ACB}=2\angle{D}$ and $x$ is the radian measure of $\angle{A}$, then $x=$
$\textbf{(A) } \frac{3\pi}{7}\qquad \textbf{(B) } \frac{4\pi}{9}\qquad \textbf{(C) } \frac{5\pi}{11}\qquad \textbf{(D) } \frac{6\pi}{13}\qquad \textbf{(E) } \frac{7\pi}{15}$

 \textbf{Answer: }A



\textbf{Solution}

We can make two equations (assume angle D is y): $y+2x=180$ and $4y+x=180$. We find that $x=\dfrac{540}{7}$. Now we have to convert this to radians. 360 degrees is $2\pi$ radians, so since we have $\dfrac{540}{7}$ degrees, the answer is $\dfrac{3\pi}{7}$ which is $\fbox{A}.$



	\item (\textit{AHSME 1990}) Four whole numbers, when added three at a time, give the sums $180,197,208$ and $222$. What is the largest of the four numbers?


 $\textbf{(A) } 77\qquad \textbf{(B) } 83\qquad \textbf{(C) } 89\qquad \textbf{(D) } 95\qquad \textbf{(E) cannot be determined from the given information}$

 \textbf{Answer: }C



\textbf{Solution}

Let $S$ be the sum of the four numbers. Then $180+197+208+222=807=3S\therefore S=269$, and subtracting $180$ gives $89$ for $\fbox{C}$



	\item (\textit{AHSME 1990}) At one of George Washington's parties, each man shook hands with everyone except his spouse, and no handshakes took place between women. If $13$ married couples attended, how many handshakes were there among these $26$ people?


 $\textbf{(A) } 78\qquad \textbf{(B) } 185\qquad \textbf{(C) } 234\qquad \textbf{(D) } 312\qquad \textbf{(E) } 325$

 \textbf{Answer: }C



\textbf{Solution}

We split this problem into two cases: A) The number of ways that men can shake hands with other men B) The number of ways that the men can shake hands with the other women (excluding their spouse).

 A) Since there are $13$ men, the number of handshakes between only men is $\frac{13 \cdot 12}{2}=78$.

 B) Since there are $13$ men and $12$ women (excluding each man's spouse), there are $13 \cdot 12 = 156$ ways.

 Adding this up, we get $156+78=234$, so the answer is$\fbox{C}$




\textbf{Solution 2}

We approach this by complement theory.

 
As there are 26 people, total handshakes between then are (26.25)/2=325

 


 Now we exclude the cases when A)handshakes of only women takes place ,

 and b) when husbands do handshakes with their spouse.

 A)(13.12)/2 =78

 B)13

 subtracting these two cases from 325, we get 234, so the answer is 234

 


	\item (\textit{AHSME 1990}) How many of the numbers, $100,101,\cdots,999$ have three different digits in increasing order or in decreasing order?


 $\textbf{(A) } 120\qquad \textbf{(B) } 168\qquad \textbf{(C) } 204\qquad \textbf{(D) } 216\qquad \textbf{(E) } 240$

 \textbf{Answer: }C



\textbf{Solution}

For decreasing order, we just need to choose any three digits from the ten: $\tbinom{10}3$

 For increasing order, the number cannot start with $0$, so choose from nine: $\tbinom93$

 The sum is $204$, giving $\fbox{C}$



	\item (\textit{AHSME 1990}) First $a$ is chosen at random from the set $\{1,2,3,\cdots,99,100\}$, and then $b$ is chosen at random from the same set. The probability that the integer $3^a+7^b$ has units digit $8$ is


 $\textbf{(A) } \frac{1}{16}\qquad \textbf{(B) } \frac{1}{8}\qquad \textbf{(C) } \frac{3}{16}\qquad \textbf{(D) } \frac{1}{5}\qquad \textbf{(E) } \frac{1}{4}$

 \textbf{Answer: }C



\textbf{Solution}

The units digits of the powers of $3$ and $7$ both cycle through $1,3,9,7$ in opposite directions, and as $4\mid 100$ each power's units digit is equally probable. There are $16$ ordered pairs of units digits, and three of them $(1,7),(7,1),(9,9)$ have a sum with units digit $8$.

 Thus the probability is $\frac3{16}$ which is $\fbox{C}$



	\item (\textit{AHSME 1990}) For how many integers $N$ between $1$ and $1990$ is the improper fraction $\frac{N^2+7}{N+4}$ $\underline{not}$ in lowest terms?


 $\textbf{(A) } 0\qquad \textbf{(B) } 86\qquad \textbf{(C) } 90\qquad \textbf{(D) } 104\qquad \textbf{(E) } 105$

 \textbf{Answer: }B



\textbf{Solution 1}

Dividing the polynomials yields $\frac{N^2 + 7}{N+4} = N - 4 + \frac{23}{N+4}$.  Since $23$ is prime, $\frac{N^2+7}{N+4}$ is not in lowest terms if and only if $N+4$ is divisible by $23$.  

 Since $N$ ranges from $1$ to $1990$, $N+4$ ranges from $5$ to $1994$.  There are $\left\lfloor\frac{1994}{23}\right\rfloor = 86$ multiples of $23$ which are at most $1994$, and thus there are $\boxed{\textbf{(B)}\ 86}$ values of $N$ such that $\frac{N^2+7}{N+4}$ is not in lowest terms.

 



\textbf{Solution 2}

For this fraction to not be in lowest terms, there must exist an integer $M > 1$ which divides both $N^2 + 7$ and $N + 4$.  Since $M$ divides $N + 4$, $N \equiv -4\ (\textrm{mod }M)$.

 Since $M$ divides $N^2+7$, $N^2 + 7 \equiv (-4)^2 + 7 \equiv 23 \equiv 0\ (\textrm{mod }M)$.  So $M$ divides $23$ and thus $M = 23$.  Therefore, this fraction is not in lowest terms if and only if $N \equiv -4\ (\textrm{mod }23)$.  That is, $N = 23k - 4$ for some integer $k$.

 The smallest such value of $N$ is $19$ (if $k = 1$), while the largest such value of $N$ is $1974$ (if $k = 86$).  Therefore, there are $\boxed{\textbf{(B)}\ 86}$ such values of $N$.

 


	\item (\textit{AHSME 1990}) 
                    \begin{figure}[H]
                        \centering
                        \includegraphics[width=0.5\textwidth]{Problem/AHSME_I_1990_20_Problem_Figure_1.png}
                    \end{figure}
                    


 In the figure $ABCD$ is a quadrilateral with right angles at $A$ and $C$. Points $E$ and $F$ are on $\overline{AC}$, and $\overline{DE}$ and $\overline{BF}$ are perpendicular to $\overline{AC}$. If $AE=3, DE=5,$ and $CE=7$, then $BF=$
$\textbf{(A) } 3.6\qquad \textbf{(B) } 4\qquad \textbf{(C) } 4.2\qquad \textbf{(D) } 4.5\qquad \textbf{(E) } 5$

 \textbf{Answer: }C



\textbf{Solution}

Label the angles as shown in the diagram. Since $\angle DEC$ forms a linear pair with $\angle DEA$, $\angle DEA$ is a right angle.

 
                    \begin{figure}[H]
                        \centering
                        \includegraphics[width=0.5\textwidth]{Solution/AHSME_I_1990_20_Solution_Figure_1.png}
                    \end{figure}
                    

 Let $\angle DAE = \alpha$ and $\angle ADE = \beta$.

 Since $\alpha + \beta = 90^\circ$, and $\alpha + \angle BAF = 90^\circ$, then $\beta = \angle BAF$. By the same logic, $\angle ABF = \alpha$.

 
As a result, $\triangle AED \sim \triangle BFA$. By the same logic, $\triangle CFB \sim \triangle DEC$.

 Then, $\frac{BF}{AF} = \frac{3}{5}$, and $\frac{CF}{BF} = \frac{5}{7}$. 

 Then, $7CF = 5BF$, and $5BF = 3AF$.

 By the transitive property, $7CF = 3AF$. $AC = AF + CF = 10$, and plugging in, we get $CF = 3$.

 Finally, plugging in to $\frac{CF}{BF} = \frac{5}{7}$, we get $BF = 4.2$ $\fbox{C}$



	\item (\textit{AHSME 1990}) Consider a pyramid $P-ABCD$ whose base $ABCD$ is square and whose vertex $P$ is equidistant from $A,B,C$ and $D$. If $AB=1$ and $\angle{APB}=2\theta$, then the volume of the pyramid is


 $\textbf{(A) } \frac{\sin(\theta)}{6}\qquad \textbf{(B) } \frac{\cot(\theta)}{6}\qquad \textbf{(C) } \frac{1}{6\sin(\theta)}\qquad \textbf{(D) } \frac{1-\sin(2\theta)}{6}\qquad \textbf{(E) } \frac{\sqrt{\cos(2\theta)}}{6\sin(\theta)}$

 \textbf{Answer: }E



\textbf{Solution}

As the base has area $1$, the volume will be one third of the height. Drop a line from $P$ to $AB$, bisecting it at $Q$.
                    \begin{figure}[H]
                        \centering
                        \includegraphics[width=0.5\textwidth]{Solution/AHSME_I_1990_21_Solution_Figure_1.png}
                    \end{figure}
                    Then $\angle QPB=\theta$, so $\cot\theta=\frac{PQ}{BQ}=2PQ$. Therefore $PQ=\frac12\cot\theta$.

 Now turning to the dotted triangle, by Pythagoras, the square of the pyramid's height is 
 \[PQ^2-(\frac12)^2=\frac{\cos^2\theta}{4\sin^2\theta}-\frac14=\frac{\cos^2\theta-\sin^2\theta}{4\sin^2\theta}=\frac{\cos 2\theta}{4\sin^2\theta}\] and after taking the square root and dividing by three, the result is $\fbox{E}$



	\item (\textit{AHSME 1990}) If the six solutions of $x^6=-64$ are written in the form $a+bi$, where $a$ and $b$ are real, then the product of those solutions with $a>0$ is


 $\textbf{(A) } -2\qquad \textbf{(B) } 0\qquad \textbf{(C) } 2i\qquad \textbf{(D) } 4\qquad \textbf{(E) } 16$

 \textbf{Answer: }D



\textbf{Solution}

This equation is $r^6e^{6\theta i}=2^6e^{(\pi\pm 2k\pi) i}$. Solving in the usual way, $r=2$ and $\theta\in\{\pm 30^\circ,\pm 90^\circ,\pm 150^\circ\}$.

 Thus there are only two solutions with positive real part, and they are conjugates, so their product is $r^2=4$ which is $\fbox{D}$



	\item (\textit{AHSME 1990}) If $x,y>0, \log_y(x)+\log_x(y)=\frac{10}{3} \text{ and } xy=144,\text{ then }\frac{x+y}{2}=$
$\textbf{(A) } 12\sqrt{2}\qquad \textbf{(B) } 13\sqrt{3}\qquad \textbf{(C) } 24\qquad \textbf{(D) } 30\qquad \textbf{(E) } 36$

 \textbf{Answer: }B



\textbf{Solution}

Let $u = \log_y(x)$.  Then $\log_x(y) = \frac{1}{u}$, so $u + \frac{1}{u} = \frac{u^2+1}{u} = \frac{10}{3}$.  

 Cross-multiplying and factoring yields $3u^2 - 10u + 3 = (3u-1)(u-3) = 0$, so $u = \frac{1}{3}$ or $u = 3$.  Thus either $x = y^3$ or $x = y^{\frac{1}{3}}$, that is, $x = y^3$ or $y = x^3$.    

 Without loss of generality, suppose $x = y^3$.  Then $y^4 = 144$ and $y = \sqrt[4]{144} = 2\sqrt{3}$.  Thus $x = 24\sqrt{3}$, so $\frac{x + y}{2} = \boxed{\textbf{(B)}\ 13\sqrt{3}}$.

 


	\item (\textit{AHSME 1990}) All students at Adams High School and at Baker High School take a certain exam. The average scores for boys, for girls, and for boys and girls combined, at Adams HS and Baker HS are shown in the table, as is the average for boys at the two schools combined. What is the average score for the girls at the two schools combined?

 \[\begin{tabular}[t]{|c|c|c|c|} \multicolumn{4}{c}{Average Scores}\\\hline Category&Adams&Baker&Adams\&Baker\\\hline Boys&71&81&79\\ Girls&76&90&?\\ Boys\&Girls&74&84& \\\hline \end{tabular}\]
$\textbf{(A) } 81\qquad \textbf{(B) } 82\qquad \textbf{(C) } 83\qquad \textbf{(D) } 84\qquad \textbf{(E) } 85$

 \textbf{Answer: }D



\textbf{Solution}

Let the numbers of boys and girls at Adams be $(A,a)$ and the numbers of boys and girls at Baker be $(B,b)$.

 Then, reading down, we have $71A+76a=74(A+a)$ and $81B+90b=84(B+b)$.Reading across, we have $71A+81B=79(A+B)$.

 Simplifying each in turn, $3A-2a=0 \implies A=2a/3$ and $3B-6b=0 \implies B=2b$ and $8A-2B=0 \implies 4A=B$.

 Thus $3b = 4a$, which implies that Adams has $\frac37$ of the girls and Baker has $\frac47$ of the girls.

 So the average score for all girls must be $\frac{3\cdot 76+4\cdot 90}7 = \frac{228+360}7 = 84$ which is $\fbox{D}$



	\item (\textit{AHSME 1990}) Nine congruent spheres are packed inside a unit cube in such a way that one of them has its center at the center of the cube and each of the others is tangent to the center sphere and to three faces of the cube. What is the radius of each sphere?


 $\textbf{(A) } 1-\frac{\sqrt{3}}{2}\qquad \textbf{(B) } \frac{2\sqrt{3}-3}{2}\qquad \textbf{(C) } \frac{\sqrt{2}}{6}\qquad \textbf{(D) } \frac{1}{4}\qquad \textbf{(E) } \frac{\sqrt{3}(2-\sqrt{2})}{4}$

 \textbf{Answer: }B



\textbf{Solution}

Let $r$ be the radius, let $C$ be the center of the cube, and let $P$ be the center of one of the eight outer spheres, noting that $PC=2r$.

 Back, in the corner of the unit cube, a smaller cube whose inner corner coincides with $P$. The cube is of dimensions $1\times 1\times 1$ and its space diagonal is of length $\sqrt3$. Thus the space diagonal of the unit cube is $2(r\sqrt3+2r)=\sqrt3$. Solving for $r$, we get $\fbox{B}$
                    \begin{figure}[H]
                        \centering
                        \includegraphics[width=0.5\textwidth]{Solution/AHSME_I_1990_25_Solution_Figure_1.png}
                    \end{figure}
                    



	\item (\textit{AHSME 1990}) Ten people form a circle.  Each picks a number and tells it to the two neighbors adjacent to them in the circle.  Then each person computes and announces the average of the numbers of their two neighbors. The figure shows the average announced by each person (\textit{not} the original number the person picked.)

                    \begin{figure}[H]
                        \centering
                        \includegraphics[width=0.5\textwidth]{Problem/AHSME_I_1990_26_Problem_Figure_1.png}
                    \end{figure}
                    
The number picked by the person who announced the average $6$ was


 $\textbf{(A) } 1 \qquad  \textbf{(B) } 5 \qquad  \textbf{(C) } 6 \qquad  \textbf{(D) } 10 \qquad \textbf{(E) }\text{not uniquely determined from the given information}$

 \textbf{Answer: }A



\textbf{Solution 1 (Ten Variables)}

For $i\in\{1,2,3,\ldots,10\},$ suppose Person $i$ picks the number $a_i$ and announces the number $i.$ We wish to find $a_6.$

 Taking the indices modulo $10,$ we are given that $\frac{a_{i-1}+a_{i+1}}{2}=i,$ from which $a_{i-1}+a_{i+1}=2i.$

 We have ten equations: five with odd-numbered indices and five with even-numbered indices. Note that these two sets of equations are independent. The set that involves $a_6$ is
 \begin{align*} a_2 + a_4 & = 6, &&(1) \\ a_4 + a_6 & = 10, &&(2) \\ a_6 + a_8 & = 14, &&(3) \\ a_8 + a_{10} & = 18, &&(4) \\ a_{10} + a_2 & = 2. &&(5) \end{align*}
Summing these five equations, we get $2(a_2 + a_4 + a_6 + a_8 + a_{10}) = 50,$ from which 
 \[a_2 + a_4 + a_6 + a_8 + a_{10} = 25. \hspace{15mm} (\bigstar)\]Subtracting $(1)+(4)$ from $(\bigstar),$ we obtain $a_6=\boxed{\textbf{(A) } 1}.$




\textbf{Solution 2 (One Variable)}

For $i\in\{1,2,3,\ldots,10\},$ suppose Person $i$ announces the number $i.$

 Let $x$ be the number picked by Person $6.$ We construct the following table:
 \[\begin{array}{c|c|c||l} & & & \\ [-2.5ex] \textbf{People} & \textbf{Average of \#s Picked} & \textbf{Sum of \#s Picked} & \multicolumn{1}{c}{\textbf{Conclusion}} \\ [0.5ex] \hline & & & \\ [-2ex] 6\text{ and }8 & 7 & 14 & \text{Person 8 picks } 14-x. \\ 8\text{ and }10 & 9 & 18 & \text{Person 10 picks } 4+x. \\ 10\text{ and }2 & 1 & 2 & \text{Person 2 picks } -2-x \\ 2\text{ and }4 & 3 & 6 & \text{Person 4 picks } 8+x \\ 4\text{ and }6 & 5 & 10 & \text{Person 6 picks } 2-x \\ \end{array}\]We have $x=2-x,$ from which $x=\boxed{\textbf{(A) } 1}.$



	\item (\textit{AHSME 1990}) Which of these triples could $\underline{not}$ be the lengths of the three altitudes of a triangle?


 $\textbf{(A) } 1,\sqrt{3},2\qquad \textbf{(B) } 3,4,5\qquad \textbf{(C) } 5,12,13\qquad \textbf{(D) } 7,8,\sqrt{113}\qquad \textbf{(E) } 8,15,17$

 \textbf{Answer: }C



\textbf{Solution}

Let $a$, $b$, and $c$ be the side lengths of the triangle such that $a<b<c$. We are given $a+b>c$ by the triangle inequality.

 Let $h_a$, $h_b$, and $h_c$ be the altitudes to sides $a$, $b$, and $c$ respectively. We see that $h_c<h_b<h_a$. By computing the areas using $a$, $b$, and $c$ as bases we get 
 \[\frac{1}{2}ah_a=\frac{1}{2}bh_b=\frac{1}{2}ch_c\]Solving for $a$ and $b$, plugging back into the triangle inequality, and canceling $c$ from both sides leaves us with 
 \[\frac{h_c}{h_a}+\frac{h_c}{h_b}>1\]Further manipulation gives 
 \[h_c>\frac{h_ah_b}{h_a+h_b}\]

 Looking at the answer choices and letting $h_c$ be the smallest value each time, we see that $5>\frac{12\cdot13}{12+13}$ is not true. Thus, the answer is $\fbox{C}$.

 


	\item (\textit{AHSME 1990}) A quadrilateral that has consecutive sides of lengths $70,90,130$ and $110$ is inscribed in a circle and also has a circle inscribed in it. The point of tangency of the inscribed circle to the side of length 130 divides that side into segments of length $x$ and $y$. Find $|x-y|$.


 $\textbf{(A) } 12\qquad \textbf{(B) } 13\qquad \textbf{(C) } 14\qquad \textbf{(D) } 15\qquad \textbf{(E) } 16$

 \textbf{Answer: }B



\textbf{Solution 1}

Let $A$, $B$, $C$, and $D$ be the vertices of this quadrilateral such that $AB=70$, $BC=110$, $CD=130$, and $DA=90$. Let $O$ be the center of the incircle. Draw in the radii from the center of the incircle to the points of tangency. Let these points of tangency $X$, $Y$, $Z$, and $W$ be on $AB$, $BC$, $CD$, and $DA$, respectively. Using the right angles and the fact that the $ABCD$ is cyclic, we see that quadrilaterals $AXOW$ and $OYCZ$ are similar.  

 Let $CZ$ have length $n$. Chasing lengths, we find that $AX=AW=n-40$. Using Brahmagupta's Formula we find that $ABCD$ has area $K=300\sqrt{1001}$ and from that we find, using that fact that $rs=K$, where $r$ is the inradius and $s$ is the semiperimeter, $r=\frac{3}{2}\sqrt{1001}$.

 From the similarity we have
 \[\frac{CY}{OX}=\frac{OY}{AX}\]Or, after cross multiplying and writing in terms of the variables, 
 \[n^2-40n-r^2=0\]Plugging in the value of $r$ and solving the quadratic gives $n=CZ=71.5$, and from there we compute the desired difference to get $\fbox{(B) 13}$.

 



\textbf{Solution 2}

Let $A, B, C, D$ be the vertices of this quadrilateral, with $AB = 70$, $BC = 90$, $CD = 130$, and $DA = 110$.  Let $\theta = \angle BCD$ and let $d = BD$.  

 Since $ABCD$ is cyclic, $\angle BAD = 180^\circ - \theta$ and thus $\cos \angle BAD = \cos(180^\circ - \theta) = -\cos \theta$.

 By the Law of Cosines on $\triangle BCD$ and $\triangle ABD$, $d^2 = 25000 - 23400 \cos \theta = 17000 + 15400 \cos \theta$.  Solving for $\cos \theta$ yields $\cos \theta = \frac{20}{97}$, and thus $\sin \theta = \sqrt{1 - \left(\frac{20}{97}\right)^2} = \frac{3\sqrt{1001}}{97}$.

 The area of $\triangle BCD$ is $\frac{1}{2}(90)(130)\sin\theta = 5850 \sin \theta$, while the area of $\triangle BAD$ is $\frac{1}{2}(70)(110)\sin(180^\circ - \theta) = 3850 \sin \theta$.  So the total area of $ABCD$ is $9700 \sin \theta = 300\sqrt{1001}$.

 Let $r$ be the inradius of $ABCD$.  Then $\frac{1}{2}r(70+90+130+110) = 300\sqrt{1001}$, so $r = \frac{3\sqrt{1001}}{2}$.

 Let $O$ be the incenter of $ABCD$, and $T$ be the foot of the perpendicular from $O$ to $CD$.  Then $\angle OCT = \frac{\theta}{2}$.  Without loss of generality, suppose $x = CT$ and $y = DT$. Then $\frac{r}{x} = \tan \frac{\theta}{2} = \frac{\sin \theta}{1 + \cos \theta} = \frac{\sqrt{1001}}{39}$.  So $x = \frac{39r}{\sqrt{1001}} = \frac{117}{2}$ and $y = 130 - \frac{117}{2} = \frac{143}{2}$.  

 Therefore, $|x-y| = \left|\frac{117}{2}-\frac{143}{2}\right| = \boxed{\textbf{(B)}\ 13}$.

 


	\item (\textit{AHSME 1990}) A subset of the integers $1,2,\cdots,100$ has the property that none of its members is 3 times another. What is the largest number of members such a subset can have?


 $\textbf{(A) } 50\qquad \textbf{(B) } 66\qquad \textbf{(C) } 67\qquad \textbf{(D) } 76\qquad \textbf{(E) } 78$

 \textbf{Answer: }D



\textbf{Solution 1}

Notice that inclusion of the integers from $34$ to $100$ is allowed as long as no integer between $12$ and $33$ inclusive is within the set. This provides a total of $100 - 34 + 1$ = 67 solutions. 

 Further analysis of the remaining integers between $1$ and $11$, we notice that we can include all the numbers except $3$ and $6$ (as including $3$ would force us to remove both $9$ and $1$ and including $6$ would make us remove $2$) to obtain the maximum number of  $9$ solutions. 

 Thus, $67 + 9 = 76$, yielding our answer, $\fbox{D}$




\textbf{Solution 2}

Write down in a column the elements $x$ which are indivisible by three, and then follow each one by $3x, 9x, 27x, \ldots$

 \[\begin{array}{ccccc}1&3&9&27&81\\2&6&18&54\\4&12&36\\5&15&45\\7&21&63\\8&24&72\\10&30&90\\11&33&99\\13&39\\\vdots&\vdots\end{array}\]We can take at most $3$ elements from the first row, and at most $2$ elements from each of the next seven rows. After that we can take only $1$ from any following row. Thus the answer is $3+7\cdot 2\,+$ the number of integers between $13$ and $100$ inclusive which are indivisible by three.

 There are $\frac13(99-12)=29$ multiples of three in that range, so there are $88-29=59$ non-multiples, and $3+14+59=76$, which is $\fbox{D}$



	\item (\textit{AHSME 1990}) If $R_n=\frac{1}{2}(a^n+b^n)$  where $a=3+2\sqrt{2}$ and $b=3-2\sqrt{2}$, and $n=0,1,2,\cdots,$ then $R_{12345}$ is an integer.  Its units digit is


 $\textbf{(A) } 1\qquad \textbf{(B) } 3\qquad \textbf{(C) } 5\qquad \textbf{(D) } 7\qquad \textbf{(E) } 9$

 \textbf{Answer: }E



\textbf{Solution 1}

$(a+b)R_n=\frac12(a^{n+1}+b^{n+1})+ab\cdot\frac12(a^{n-1}+b^{n-1})=R_{n+1}+abR_{n-1}$but $a+b=6$ and $ab=1$, so this means that $R_{n+1}=6R_n-R_{n-1}$. Since $R_0=1$ and $R_1=3$, all terms are integers and we can continue the recurrence $\textrm{mod}\ 10$ to get the repeating sequence $1,3,7,9,7,3$. The number $12345$ is divisible by three but not by six, so the required element in the sequence is $9$ which is answer $\fbox{E}$




\textbf{Solution 2 (Characteristic Polynomials)}

Notice that $3 \pm 2\sqrt{2} = 3 \pm \sqrt{8}$ must be the roots of $(x - 3)^2 - 8$, or $x^2 - 6x + 1$.  So the characteristic polynomial of $\{R_n\}$ is $x^2 - 6x + 1$, and thus the recurrence $R_n = 6R_{n-1} - R_{n-2}$ holds for all $n \geq 2$.

 Since $R_0 = 1$ and $R_1 = 3$, $R_n$ is an integer for all $n$.  Let $S_n$ be the last digit of $R_n$.  Using the recurrence relation, the first terms of $\{S_n\}$ starting with $S_0$ are $1, 3, 7, 9, 7, 3, 1, 3, ...$ 

 Since each term of this sequence depends solely on the previous two terms, this sequence cycles every $6$ terms.  The remainder when $12345$ is divided by $6$ is $3$, so $S_{12345} = S_3 = \boxed{\textbf{(E)}\ 9}$.

 -j314andrews

 



\textbf{Solution 3}

Notice that for all integers $n \geq 0$, $\left(3 + 2\sqrt{2}\right)^n$ can be written as $x + y\sqrt{2}$, where $x$ and $y$ are integers.  Furthermore, if $\left(3 + 2\sqrt{2}\right)^n = x + y\sqrt{2}$, where $x$ and $y$ are integers, then $\left(3 - 2\sqrt{2}\right)^n = x - y\sqrt{2}$.

 Let $\{x_n\}$, $\{y_n\}$ be sequences of integers such that for all integers $n \geq 0$, $\left(3 + 2\sqrt{2}\right)^n = x_n + y_n\sqrt{2}$.  Then $\left(3 + 2\sqrt{2}\right)^{n+1} = \left(3 + 2\sqrt{2}\right)\left(x_n + y_n\sqrt{2}\right) = (3x_n + 4y_n) + (2x_n + 3y_n)\sqrt{2}$.  So $x_{n+1} = 3x_n + 4y_n$ and $y_{n+1} = 2x_n + 3y_n$.  

 The first few ordered pairs $(x_n, y_n)$ modulo $10$, starting with $(x_0, y_0)$, are $(1, 0), (3, 2), (7, 2), (9, 0), (7, 8), (3, 8), (1, 0), ...$.  Since each ordered pair in this sequence depends solely on the previous ordered pair, this sequence must cycle every $6$ entries.

 Since $R_n = \frac{1}{2}\left((x_n + y_n \sqrt{2}) + (x_n - y_n \sqrt{2})\right) = x_n$, the last digit of $R_{12345}$ is $x_{12345}$ modulo $10$.  Since $12345 \equiv 3\ (\text{mod}\ 6)$, $x_{12345} \equiv x_{3} \equiv \boxed{\textbf{(E)}\ 9}\ (\text{mod}\ 10)$.

 -j314andrews

 


\end{enumerate}